% Options for packages loaded elsewhere
\PassOptionsToPackage{unicode}{hyperref}
\PassOptionsToPackage{hyphens}{url}
\PassOptionsToPackage{dvipsnames,svgnames,x11names}{xcolor}
%
\documentclass[
]{article}
\usepackage{amsmath,amssymb}
\usepackage{iftex}
\ifPDFTeX
  \usepackage[T1]{fontenc}
  \usepackage[utf8]{inputenc}
  \usepackage{textcomp} % provide euro and other symbols
\else % if luatex or xetex
  \usepackage{unicode-math} % this also loads fontspec
  \defaultfontfeatures{Scale=MatchLowercase}
  \defaultfontfeatures[\rmfamily]{Ligatures=TeX,Scale=1}
\fi
\usepackage{lmodern}
\ifPDFTeX\else
  % xetex/luatex font selection
\fi
% Use upquote if available, for straight quotes in verbatim environments
\IfFileExists{upquote.sty}{\usepackage{upquote}}{}
\IfFileExists{microtype.sty}{% use microtype if available
  \usepackage[]{microtype}
  \UseMicrotypeSet[protrusion]{basicmath} % disable protrusion for tt fonts
}{}
\makeatletter
\@ifundefined{KOMAClassName}{% if non-KOMA class
  \IfFileExists{parskip.sty}{%
    \usepackage{parskip}
  }{% else
    \setlength{\parindent}{0pt}
    \setlength{\parskip}{6pt plus 2pt minus 1pt}}
}{% if KOMA class
  \KOMAoptions{parskip=half}}
\makeatother
\usepackage{xcolor}
\setlength{\emergencystretch}{3em} % prevent overfull lines
\providecommand{\tightlist}{%
  \setlength{\itemsep}{0pt}\setlength{\parskip}{0pt}}
\setcounter{secnumdepth}{-\maxdimen} % remove section numbering
% ---------------------------------------------------------------------------
% header.tex — pandoc MD→PDF code block styling
% Pair with: --highlight-style .fcc/pdf/p10k.theme
%
% Approach: tcolorbox owns the block background entirely.
% fvextra patches Highlighting directly to fix the ␣ bug.
% ---------------------------------------------------------------------------

% ---------------------------------------------------------------------------
% Page margins
% ---------------------------------------------------------------------------
\usepackage[top=2cm, bottom=2cm, left=2cm, right=2cm]{geometry}

% ---------------------------------------------------------------------------
% Hyperlinks — controlled via -V colorlinks=true passed to pandoc.
% pandoc's template conditionally sets colorlinks/hidelinks based on that
% variable — no hyperref config needed here.
% ---------------------------------------------------------------------------

% ---------------------------------------------------------------------------
% Images — required for \includegraphics in title page template
% ---------------------------------------------------------------------------
\usepackage{graphicx}
\usepackage{xcolor}
\usepackage{fvextra}
\usepackage{etoolbox}
\usepackage[most]{tcolorbox}

% ---------------------------------------------------------------------------
% p10k colours (xterm-256 exact hex)
% ---------------------------------------------------------------------------
\definecolor{termbg} {HTML}{262626}
\definecolor{termfg} {HTML}{bcbcbc}

% ---------------------------------------------------------------------------
% Fix ␣ bug: patch the Highlighting environment that pandoc defines via
% \DefineVerbatimEnvironment{Highlighting}{Verbatim}{commandchars=...}
% fvextra's \fvset doesn't reach it because the options on DefineVerbatimEnvironment
% take precedence. We redefine it after \begin{document} to inject breaklines.
% ---------------------------------------------------------------------------
\AtBeginDocument{%
  % Redefine Highlighting to add fvextra breaklines + correct foreground colour.
  % Preserve the commandchars option pandoc requires for token coloring.
  \DefineVerbatimEnvironment{Highlighting}{Verbatim}{%
    commandchars=\\\{\},%
    breaklines=true,%
    breakanywhere=false,%
    breaksymbolleft={},%
    fontsize=\small,%
  }%
  % Suppress hyphenation at font level — prevents TeX inserting hyphens on line breaks.
  \BeforeBeginEnvironment{Verbatim}{\hyphenchar\font=-1\relax}%
  \AfterEndEnvironment{Verbatim}{\hyphenchar\font=`\-\relax}%
}

% ---------------------------------------------------------------------------
% Redefine Shaded using tcolorbox — completely independent of framed/snugshade.
% \AtBeginDocument so it fires after pandoc defines Shaded at line ~47.
% ---------------------------------------------------------------------------
\AtBeginDocument{%
  \ifdef{\Shaded}{}{\newenvironment{Shaded}{}{}}%
  \renewenvironment{Shaded}{%
    \begin{tcolorbox}[%
      enhanced,%
      arc=3pt,%
      boxrule=0.5pt,%
      colback=termbg,%
      colframe=termbg!60!black,%
      coltext=termfg,%
      left=6pt, right=6pt, top=4pt, bottom=4pt,%
      before skip=8pt, after skip=8pt,%
      fontupper=\color{termfg}\small,%
    ]%
  }{%
    \end{tcolorbox}%
  }%
}

% ---------------------------------------------------------------------------
% Inline code styling
% Use \scalerel to allow the text to shrink when inside narrow table columns.
% \texttt globally uses \footnotesize so long identifiers fit in pandoc's
% hardcoded longtable column proportions.
% ---------------------------------------------------------------------------
\usepackage{relsize}
\let\oldtexttt\texttt
\renewcommand{\texttt}[1]{%
  \colorbox{termbg}{{\ttfamily\relsize{-1}\color{termfg}#1}}%
}

% ---------------------------------------------------------------------------
% Long strings - see breakable.lua, which emits these two.
% \hcbreak: a break opportunity after a separator (/ . - _ ...) in a long
% word, URL or identifier. \hcsplit: one inside a run with no separator - a
% higher penalty, so TeX takes it only when nothing better fits. Both cost
% more than a space, so TeX still prefers to break between words.
% \hccodepart{left}{right}{text}: one piece of a long inline code span, styled
% as \texttt above. A piece that continues on a side gets \kern-\fboxsep there,
% so its box has no padding on that side; the \vphantom gives every piece the
% same height and depth, so the pieces join without a step.
% xurl: \url may break at any character.
% ---------------------------------------------------------------------------
\newcommand{\hcbreak}{\penalty50\relax}
\newcommand{\hcsplit}{\penalty300\relax}
\newcommand{\hccodepart}[3]{%
  \colorbox{termbg}{#1{\ttfamily\relsize{-1}\color{termfg}\vphantom{Xgjpqy\textunderscore}#3}#2}%
}
\usepackage{xurl}

% ---------------------------------------------------------------------------
% Suppress all hyphenation in prose.
% \pretolerance=10000 makes TeX accept slightly loose lines before trying
% hyphenation. Combined with \hyphenpenalty=10000 it effectively disables
% hyphenation without breaking the layout engine entirely.
% ---------------------------------------------------------------------------
\pretolerance=10000
\hyphenpenalty=10000
\exhyphenpenalty=10000

% ---------------------------------------------------------------------------
% Table row colours — alternating light grey / white body rows, medium grey
% header. Applies to all longtables pandoc generates.
%
% \rowcolors{<start row>}{<odd colour>}{<even colour>}
% Start row 1 so the header row is row 0 (styled separately via \rowcolor).
% \arrayrulecolor keeps the rule lines visible against the light background.
% ---------------------------------------------------------------------------
\usepackage{colortbl}
\definecolor{tableheader}{HTML}{cccccc}
\definecolor{tablerowodd}{HTML}{f5f5f5}
\definecolor{tableroweven}{HTML}{ffffff}

\AtBeginDocument{%
  \rowcolors{2}{tablerowodd}{tableroweven}%
}

% Style the header row — pandoc wraps it in \endfirsthead / \endhead blocks.
% Redefine \toprule and \midrule to inject \rowcolor on the header row.
% \arrayrulecolor darkens the rule lines so they remain visible on light bg.
\arrayrulecolor{gray!50}

% Colour the header row by patching longtable's first-head block.
% \pretocmd injects \rowcolor{tableheader} immediately after \toprule,
% which is where pandoc places the header row.
\usepackage{longtable}
\AtBeginDocument{%
  \pretocmd{\toprule}{\rowcolor{tableheader}}{}{}%
}
\ifdefined\setmonofont
  \setmonofont{DejaVu Sans Mono}
\fi
% fallback.tex - written by holo-convert for each run (fonts.py).
% Characters Latin Modern has no glyph for used to vanish from the PDF, with
% only a warning per character from pandoc. XeTeX's character classes switch
% to DejaVu for exactly those Unicode blocks and back after them. Code (the
% monospace family) is left alone: its font has these, and ASCII art there
% must keep its fixed width. No groups, so a brace between the two ends of a
% run cannot unbalance anything.
\ifdefined\XeTeXinterchartoks
\newfontfamily\hcfallbackletters{DejaVu Serif}
\newfontfamily\hcfallbacksymbols{DejaVu Sans}
\newXeTeXintercharclass\hcclassletters
\newXeTeXintercharclass\hcclasssymbols
\makeatletter
\def\hc@range#1#2#3{\count@=#1\relax
  \loop\XeTeXcharclass\count@=#3\relax\ifnum\count@<#2\advance\count@\@ne\repeat}
\hc@range{"0370}{"03FF}{\hcclassletters}
\hc@range{"1F00}{"1FFF}{\hcclassletters}
\hc@range{"0400}{"052F}{\hcclassletters}
\hc@range{"2196}{"21FF}{\hcclasssymbols}
\hc@range{"2200}{"22FF}{\hcclasssymbols}
\hc@range{"2300}{"23FF}{\hcclasssymbols}
\hc@range{"2460}{"24FF}{\hcclasssymbols}
\hc@range{"2500}{"257F}{\hcclasssymbols}
\hc@range{"2580}{"259F}{\hcclasssymbols}
\hc@range{"25A0}{"25FF}{\hcclasssymbols}
\hc@range{"2600}{"26FF}{\hcclasssymbols}
\hc@range{"2700}{"27BF}{\hcclasssymbols}
\hc@range{"27F0}{"27FF}{\hcclasssymbols}
\hc@range{"2900}{"297F}{\hcclasssymbols}
\hc@range{"2980}{"29FF}{\hcclasssymbols}
\hc@range{"2A00}{"2AFF}{\hcclasssymbols}
\hc@range{"2B00}{"2BFF}{\hcclasssymbols}
\let\hc@close\relax
\def\hc@restore{\fontfamily{\hc@saved}\selectfont\let\hc@close\relax}
\def\hc@open#1{\hc@close
  \edef\hc@a{\f@family}\edef\hc@b{\ttdefault}%
  \ifx\hc@a\hc@b\else\let\hc@saved\hc@a\let\hc@close\hc@restore#1\fi}
\XeTeXinterchartoks 0 \hcclassletters = {\hc@open\hcfallbackletters}
\XeTeXinterchartoks 4095 \hcclassletters = {\hc@open\hcfallbackletters}
\XeTeXinterchartoks \hcclasssymbols \hcclassletters = {\hc@open\hcfallbackletters}
\XeTeXinterchartoks 0 \hcclasssymbols = {\hc@open\hcfallbacksymbols}
\XeTeXinterchartoks 4095 \hcclasssymbols = {\hc@open\hcfallbacksymbols}
\XeTeXinterchartoks \hcclassletters \hcclasssymbols = {\hc@open\hcfallbacksymbols}
\XeTeXinterchartoks \hcclassletters 0 = {\hc@close}
\XeTeXinterchartoks \hcclassletters 4095 = {\hc@close}
\XeTeXinterchartoks \hcclasssymbols 0 = {\hc@close}
\XeTeXinterchartoks \hcclasssymbols 4095 = {\hc@close}
\makeatother
\XeTeXinterchartokenstate=1
\fi
\ifLuaTeX
  \usepackage{selnolig}  % disable illegal ligatures
\fi
\usepackage{bookmark}
\IfFileExists{xurl.sty}{\usepackage{xurl}}{} % add URL line breaks if available
\urlstyle{same}
\hypersetup{
  colorlinks=true,
  linkcolor={blue},
  filecolor={Maroon},
  citecolor={blue},
  urlcolor={blue},
  pdfcreator={LaTeX via pandoc}}

\author{}
\date{}

\begin{document}

\section{Characters beyond Latin}\label{characters-beyond-latin}

Arrows and marks: A ⇒ B, done ✓, not done ✗, → and ←, x ≤ y, ★.

Greek and Cyrillic: αβγ δ --- Привет, мир.

Box drawing in prose: ┌──┐ │ok│ └──┘

\begin{verbatim}
┌────────┬────────┐
│ ascii  │ art    │
└────────┴────────┘
\end{verbatim}

Not supported in PDF output, and reported: 你好 こんにちは 안녕 🚀 שלום
مرحبا

\end{document}
